%% macros.tex — paquetes y macros para el anteproyecto TFG
%% Basado en la plantilla de P. Alvarado, EIE-TEC

\usepackage{ifpdf}
\newcommand{\ifdraft}[2]{\ifthenelse{\boolean{draftmode}}{#1}{#2}}
\newcommand{\ifbook}[2]{\ifthenelse{\boolean{bookmode}}{#1}{#2}}

%% Idioma y codificación
\usepackage{csquotes}
\usepackage[spanish,es-noindentfirst]{babel}

%% Matemáticas
\usepackage{amsmath}
\usepackage{amssymb,amstext}
\usepackage{nicefrac}

%% Tablas y figuras
\usepackage{array}
\usepackage{colortbl}
\usepackage{longtable}
\usepackage{tabularx}
\usepackage{booktabs}
\usepackage{multirow}
\usepackage{multicol}
\usepackage{float}
\usepackage[format=hang,font=small,labelfont=bf]{caption}
\usepackage{graphicx}

%% Listas
\usepackage{enumitem}
\newlist{compactitem}{itemize}{3}
\setlist[compactitem]{topsep=0pt,partopsep=0pt,itemsep=0pt,parsep=0pt}
\setlist[compactitem,1]{label=\textbullet}
\setlist[compactitem,2]{label=---}
\setlist[compactitem,3]{label=*}

%% Tipografía
\usepackage[expansion=false]{microtype}

%% Estilos de página
\usepackage{vmargin}
\setpapersize{USletter}
\setmarginsrb{33mm}{8mm}{23mm}{7mm}{15pt}{15pt}{7mm}{12mm}

\usepackage{fancyhdr}
\usepackage{lastpage}
\usepackage{afterpage}

%% Interlineado
\renewcommand{\baselinestretch}{1.1}
\parindent0em
\parskip1.3ex

%% Profundidad de numeración y TOC
\setcounter{secnumdepth}{2}
\setcounter{tocdepth}{2}

%% Estilo TEC
\ifdraft{%
  \usepackage[todo]{sty/tecStyle}
}{%
  \usepackage{sty/tecStyle}
}

%% TOC bind
\usepackage[nottoc]{tocbibind}

%% URLs
\usepackage{url}
\usepackage{rotating}

%% Biblatex (cargado en main.tex — no cargar aquí de nuevo)

%% Placeholder para figuras pendientes (no requiere archivo externo)
\newcommand{\figpend}[1][0.8\textwidth]{%
  \noindent\makebox[#1]{%
    \fbox{\begin{minipage}[c][4cm]{#1}%
      \centering\color{gray}\textit{[Figura pendiente]}%
    \end{minipage}}%
  }}

%% Hyperref (siempre al final)
\ifpdf
  \ifdraft{%
    \usepackage[naturalnames=true,linktocpage,hyperindex,colorlinks,
                urlcolor=dkred,filecolor=dkmagenta,linkcolor=dkred,
                citecolor=dkgreen,plainpages=false,pdfpagelabels,
                pdfpagemode=UseOutlines,bookmarksnumbered=true,
                pdfpagelayout=OneColumn,pdfview=FitH,
                pdfstartview=FitH]{hyperref}
  }{%
    \usepackage[naturalnames=true,linktocpage,hyperindex,colorlinks,
                urlcolor=dkred,filecolor=dkmagenta,linkcolor=dkred,
                citecolor=dkgreen,plainpages=false,pdfpagelabels,
                pdfpagemode=UseOutlines,bookmarksnumbered=true,
                pdfpagelayout=OneColumn,pdfview=FitH,
                pdfstartview=FitH]{hyperref}
  }
  \usepackage[figure]{hypcap}
  \pdfadjustspacing=1
\else
  \usepackage{graphicx}
  \usepackage[ps2pdf,linktocpage,hyperindex,
              pdfpagemode=UseOutlines,pdfstartview=FitH]{hyperref}
\fi

\AtBeginDocument{%
  \hypersetup{pdftitle={\pdfTitle},pdfsubject={\pdfSubject},
              pdfauthor={\pdfAuthor},pdfkeywords={\pdfKeywords}}%
}

%% Headers y footers
\newcommand{\copyrightfooter}{\tiny{\copyright \the\year\ --- \scriptAuthorShort \qquad Uso exclusivo ITCR}}
\newcommand{\draftfoot}{\ifdraft{\textcolor{dkblue}{\tiny\textsl{Borrador: \today}}}{}}

\pagestyle{fancy}
\renewcommand{\chaptermark}[1]{\markboth{{\small\thechapter\hspace*{1mm}#1}}{}}
\renewcommand{\sectionmark}[1]{\markright{{\small\thesection\hspace*{1mm}#1}}{}}
\lhead[{\small\thepage}]{\fancyplain{}{{\slshape\small\nouppercase{\leftmark}}}}
\chead[]{}
\rhead[\fancyplain{}{{\slshape\small\nouppercase{\rightmark}}}]{{\small\thepage}}
\lfoot[]{\draftfoot}
\ifbook{\cfoot[]{}}{\cfoot[\copyrightfooter]{\copyrightfooter}}
\rfoot[\draftfoot]{}
\renewcommand{\headrulewidth}{0.5pt}
\renewcommand{\footrulewidth}{0pt}

%% Caption para tablas
\captionsetup[table]{position=top,format=hang,
  textfont={normalsize},labelfont={normalsize,bf}}
\addto\extrasspanish{\renewcommand{\tablename}{Tabla}}
\addto\extrasspanish{\renewcommand{\listtablename}{\'Indice de tablas}}

%% Comandos para tribunal
\newcommand{\nameLectorI}{$<$\emph{Use setLectorI en main.tex}$>$}
\newcommand{\genderLectorI}{$<$\emph{Use setLectorI en main.tex}$>$}
\newcommand{\setLectorI}[1][M]{%
  \ifthenelse{\equal{#1}{F}}{\renewcommand{\genderLectorI}{Profesora Lectora}}%
                             {\renewcommand{\genderLectorI}{Profesor Lector}}
  \lectorIRelay}
\newcommand{\lectorIRelay}[1]{\renewcommand{\nameLectorI}{#1}}

\newcommand{\nameLectorII}{$<$\emph{Use setLectorII en main.tex}$>$}
\newcommand{\genderLectorII}{$<$\emph{Use setLectorII en main.tex}$>$}
\newcommand{\setLectorII}[1][M]{%
  \ifthenelse{\equal{#1}{F}}{\renewcommand{\genderLectorII}{Profesora Lectora}}%
                             {\renewcommand{\genderLectorII}{Profesor Lector}}
  \lectorIIRelay}
\newcommand{\lectorIIRelay}[1]{\renewcommand{\nameLectorII}{#1}}

\newcommand{\nameAsesor}{$<$\emph{Use setAsesor en main.tex}$>$}
\newcommand{\genderAsesor}{Profesor Asesor}
\newcommand{\setAsesor}[1][M]{%
  \ifthenelse{\equal{#1}{F}}{\renewcommand{\genderAsesor}{Profesora Asesora}}%
                             {\renewcommand{\genderAsesor}{Profesor Asesor}}
  \asesorRelay}
\newcommand{\asesorRelay}[1]{\renewcommand{\nameAsesor}{#1}}
